\section{The Commission and the Judgment of 1862}

\subsection{The commission of 28 July 1858}

Bishop Laurence moved quickly by the standards of episcopal
prudence: the eighteenth reported apparition was on 16 July 1858,
and on 28 July he constituted a commission of inquiry. The
Sanctuary transmits the commission's mandate from the episcopal
ordinance: ``to collect and to ascertain the facts which have
happened or which could still occur in or around the grotto of
Lourdes, to report them to us, to highlight their nature, and thus
to provide us with the components necessary for us to reach a
conclusion''; and, in the French of the Sanctuary's anniversary
notice, the commission was ``charg\'ee de constater
l'authenticit\'e et la nature des faits qui se sont
produits.''\endnote{Sanctuary of Lourdes, ``The Decree,''
\url{https://www.lourdes-france.org/en/the-decree/}, and
``Recognition of the apparitions,''
\url{https://www.lourdes-france.com/en/recognition-of-the-apparitions/};
French mandate clause from ``1862--2022: 160 ans de la
reconnaissance des Apparitions,''
\url{https://www.lourdes-france.com/1862-2022-160-ans-de-la-reconnaissance-des-apparitions-de-lourdes/},
all read 2026-07-25. The 160th-anniversary notice also records
that the Sanctuary prepared a special re-edition of the mandement
in 2022; that re-edition was not consulted.}

The questions the commission pursued, as the Sanctuary summarizes
its work, were the sober ones: had cures in fact occurred after
use of the grotto water; was the water natural or supernatural;
were Bernadette's visions real, and if real, did they have a
divine character; had the apparition made requests, and which; had
the spring existed before 25 February. The bishop's ordinance
insisted on breadth of method---establishing facts, questioning
witnesses, consulting men of science, ``in particular the doctors
who would have treated the sick before their recovery, but also
experts in the fields of physics, chemistry, geology''---and on an
open outcome: ``The commission must not neglect a single thing to
obtain all the information possible in order to arrive at the
truth, whatever it may be.''\endnote{``Recognition of the
apparitions'' page, quoting the ordinance in its own English; the
page is the Sanctuary's institutional presentation of its founding
documents, and the quotations are used at that ceiling
(Appendix~A). The commission's composition---``pious, learned, and
experienced priests''---is the mandement's own description
(Lasserre 1870, p.~400).} The commission interrogated Bernadette
itself---the mandement speaks of hearing her ``in presence of the
Commission, assembled to hear her a second time''---and worked for
nearly four years before the bishop concluded.\endnote{Lasserre
1870, p.~399 (``en pr\'esence de la Commission, r\'eunie pour
l'entendre une seconde fois''); p.~400 (``L'\'ev\'enement dont
nous vous entretenons est, depuis quatre ann\'ees, l'objet de
notre sollicitude'').}

\subsection{The mandement's argumentation}

The mandement of 18 January 1862 is not a bare decree; it argues
its conclusion across a doctrinal preamble, a narrative, a
discussion, and formal considerations, and the argument deserves
to be reported in its own structure, because its restraint is part
of what the Church later approved of in it.

It opens by setting apparitions within salvation history and then
states the Church's method in a sentence that could serve as the
epigraph of the whole Lourdes dossier: ``the Church brings a wise
slowness to the appreciation of supernatural facts: she demands
certain proofs before admitting them and proclaiming them
divine,'' since the devil ``sometimes transforms himself into an
angel of light''; the bishop cites the Johannine command to test
the spirits.\endnote{Lasserre 1870, pp.~399--400: ``l'\'Eglise
apporte une sage lenteur dans l'appr\'eciation des faits
surnaturels: elle demande des preuves certaines, avant de les
admettre et de les proclamer divins,'' with 2~Cor 11:14 and
1~Jn 4:1 in the mandement's own notes.} The test, he says, has
been made: ``We have kept ourselves informed through the
Commission, composed of pious, learned, and experienced priests,
who questioned the child, studied the facts, examined everything,
weighed everything. We have also invoked the authority of Science,
and we have remained convinced that the Apparition is supernatural
and divine, and that, consequently, what Bernadette saw is the
Most Holy Virgin. Our conviction was formed on the testimony of
Bernadette, but above all on the facts that have occurred, which
can be explained only by a divine intervention.''\endnote{Lasserre
1870, p.~400 (editorial translation; French quoted in the note at
\S2.2 and in the source audit). The Sanctuary's English
``Recognition'' page renders the same passage; the two renderings
were compared and agree in substance.}

The discussion of Bernadette's testimony then runs on two rails.
Sincerity: her simplicity and modesty, her silence unless
questioned, her constancy across interrogations ``without adding
anything, without retracting anything,'' her imperviousness to
threats and to offers---``the sincerity of Bernadette is therefore
incontestable. Let us add that it is uncontested,'' even by her
contradictors. Sanity: the wisdom of her answers, a calm
imagination, no intellectual disorder, no morbid affection; and
then the argument this study flagged at \S2.3---she saw not once
but eighteen times, saw suddenly the first time when nothing
prepared her, and \emph{failed} to see on two days of the
fortnight though the circumstances were identical; therefore
hallucination is excluded, ``the young girl really saw and heard a
being calling herself the Immaculate Conception,'' and since the
phenomenon cannot be explained naturally, ``we are justified in
believing that the Apparition is
supernatural.''\endnote{Lasserre 1870, pp.~400--402. The
italicized failure-to-see is the mandement's own evidential use of
the empty days, an argument from the phenomenon's independence of
the visionary's expectation.}

But the bishop does not rest the judgment on the visionary alone.
The testimony ``borrows a wholly new force---we shall even say its
complement---from the marvelous facts'' that followed: the
immediate, immense, and persevering concourse of the faithful; the
conversions and returns to practice attributed to the invocation
of Notre-Dame de Lourdes---``marvels of grace, which bear a
character of universality and duration, [and] can have only God
for author''---and the cures. On the cures the mandement is
specific in a way that anticipates the whole later medical
apparatus: the sick used water ``deprived of any natural curative
quality, on the report of skilled chemists who made a rigorous
analysis of it''; cures came ``some instantaneously, others after
two or three uses''; the cured included the ``abandoned and
declared incurable''; the cures ``are permanent''; asked whether
organic power could have produced them, ``Science, consulted on
this subject, answered negatively.'' ``These cures are therefore
the work of God.''\endnote{Lasserre 1870, pp.~402--404. The
sequence---serious illness, no natural curative agent,
instantaneous or near-instantaneous cure, permanence, negative
scientific verdict---is materially the profile the Lambertini
criteria formalize and the modern process still uses (\S5.3).}
The peroration claims the classical sign: \emph{Digitus Dei est
hic}---and immediately reads the event against the dogma of 1854
(\S3.3).

The formal considerations before the operative articles bind the
reasoning to authority: the bishop grounds himself ``on the rules
wisely traced by Benedict XIV in his work on the Beatification and
Canonization of saints, for the discernment of true and false
apparitions''; cites the favorable report of the commission and
``the written testimony of the Doctors-Physicians whom we
consulted about numerous cures''; and reasons in three
considerations that the fact ``cannot be explained except through
a supernatural cause,'' that the cause ``can only be divine''
since the effects are on one side sensible signs of grace (the
conversion of sinners) and on the other ``derogations from the
laws of nature'' (the miraculous cures), and that the conviction
is fortified by the immense and spontaneous concourse of the
faithful.\endnote{Lasserre 1870, pp.~406--407, including the
mandement's own note citing Benedict XIV, \work{De servorum Dei
beatificatione et beatorum canonizatione}, book~III (the printed
note reads ``Liv. III, ch. ii''; the exact chapter reference was
not verified against an edition of Benedict XIV and is reported as
printed).}

\subsection{The operative articles and their exact object}

The operative articles are printed and translated at \S1; what
needs analysis here is their scope, clause by clause.

\begin{statusdossier}{What the 1862 act judged---and what it did
not}
\dosfield{Event judged}{The apparition of ``the Immaculate Mary,
Mother of God'' to Bernadette Soubirous, ``on 11 February 1858 and
following days, eighteen times,'' at the grotto of Massabielle:
the bounded eighteen-event corpus of \S2.3, no more.}
\dosfield{Modality}{``The faithful are justified in believing it
certain'' (\term{fond\'es \`a la croire certaine}): a judgment of
credibility warranting assent, in the register of the
pre-conciliar discernment of apparitions; it neither imposes
belief nor converts the event into public Revelation (\S8).}
\dosfield{Cult}{Authorized ``in our diocese''---a diocesan
authorization of the cult of Notre-Dame de la Grotte de Lourdes;
universal liturgical extension came only by later Roman acts
(\S7.1). Devotional publications remained under written episcopal
approval---a control clause, not a blank endorsement of Lourdes
literature.}
\dosfield{Cures}{Seven cures were recognized as miraculous in
connection with the same act, on the examination of Professor
Vergez and the consulted physicians (\S4.4); the act did not
canonize any general theory of the water, which its own
argumentation had described as naturally inert.}
\dosfield{Submission}{The judgment is expressly submitted ``to the
judgment of the Sovereign Pontiff''; no papal act of confirmation
is claimed by the mandement, and none is asserted by this study.
Later papal engagement (\S7) received and honored the judgment
without re-issuing it.}
\dosfield{Not judged}{The verbal inerrancy of any reported
sentence; the candle episode of 7 April as a distinct prodigy
(\S5.5); the private prayer and secrets; any later narrative,
image, book, or devotion (expressly reserved); any physical
mechanism; any obligation of divine faith.}
\dosfield{Witness ceiling}{Operative wording quoted from Lasserre's
1870 reproduction, collated with the Sanctuary's French and
English transcriptions (readings agree); the 1862 diocesan imprint
and archive original were not consulted (Appendix~A).}
\end{statusdossier}

The signature and date clause close the act in the classical
diocesan form: given at Tarbes, under the bishop's sign and seal
and the counter-signature of the chancery secretary (Fourcade,
canon), ``le 18 janvier 1862, f\^ete de la Chaire de Saint-Pierre
\`a Rome.''\endnote{Lasserre 1870, p.~409. The choice of the
feast of the Chair of Peter harmonizes with the submission clause
that precedes it; the mandement was ordered read in all churches,
chapels, and oratories of the diocese on the Sunday after
reception.}

\subsection{The seven cures of 18 January 1862}

The same act carried the first cure recognitions. The Sanctuary's
institutional history states that the bishop ``used three criteria
to establish that the Immaculate Conception had actually appeared
to Bernadette Soubirous: reliability of the seer, spiritual
fruits, cures of the body,'' and that on 18 January 1862 ``he
declared by Decree that seven cures, among the hundreds examined
by Prof Henri Vergez, were `miraculous.'\,''\endnote{Sanctuary of
Lourdes, ``The Medical Bureau of the Sanctuary,''
\url{https://www.lourdes-france.org/en/medical-bureau-sanctuary/},
read 2026-07-25. Vergez is identified here only as the page
identifies him (``Prof Henri Vergez''); his appointment and
faculty were not separately verified.} The
Sanctuary's list of recognized cures names the seven, each with
recognition date 18 January 1862: Catherine Latapie (Loubajac),
Louis Bouriette (Lourdes), Blaisette Cazenave (Lourdes), Henri
Busquet (Nay), Justin Bouhort (Lourdes), Madeleine Rizan (Nay),
and Marie Moreau (Tartas).\endnote{Sanctuary of Lourdes,
``Miraculous healings,''
\url{https://www.lourdes-france.com/en/miraculous-healings/},
entries 1--7. The list gives recognition dates; the cure events
themselves belong to 1858 in the received narrative (Latapie and
Bouriette at \S2.6), and individual cure dates are not asserted
here beyond it.}

Two things follow from the 1862 recognitions that shape everything
in \S5. First, the juridical form was set at the beginning: a
medical examination (here Vergez and the consulted physicians)
feeding an episcopal declaration---the physician never declares
the miracle, the bishop never certifies the medicine. Second, the
evidential standard was already the demanding one: the mandement's
own description of the cures it accepted---grave illness, no
natural curative agent, suddenness, permanence, negative verdict
of consulted science---is recognizably the profile that the
Lambertini criteria would supply to the modern process. The
Lourdes medical tradition did not soften the standard over time;
it institutionalized it.
